\clearpage\begingroup\SourceLeaf{682b}
\ArithmeticHead{47}{R}{IV. POTENTIALTHEORIE.}
\makeatletter\tagsleft@true\let\veqno\@@leqno\makeatother
\renewcommand{\thefootnote}{\arabic{footnote})}\setcounter{footnote}{85}
\begin{center}\textls[100]{2. Die Entwicklung bei \textsc{Gauss}.}\end{center}
\noindent\hspace*{\parindent}Wir wenden uns nun zur historischen Entwicklung der Potentialtheorie
bei \textsc{Gauss}. Seine Arbeiten auf diesem Gebiete entspringen aus drei Quellen:
1. der Beschäftigung mit dem Attraktionsproblem der Ellipsoide; 2. den Studien
über Funktionentheorie und konforme Abbildung; 3. den Arbeiten zur Theorie
des Erdmagnetismus.

1. Die \textit{Theoria attractionis} von 1813 enthält vor allem die sechs Integralsätze
zur Reduktion dreifacher Integrale auf Oberflächenintegrale und die
entsprechenden Verfahren der Zerlegung in prismatische und kegelförmige
Elemente; diese Sätze sind spätestens 1812 entstanden. Im Anschluss an
sie führt eine Werke XI, 1, S. 71 abgedruckte Nachlassnotiz aus Handbuch
21, Bg, die auf 1814 zu datieren ist, den Begriff der Kraftröhre und der
dazu orthogonalen Niveauflächen, sowie der dadurch bedingten Intensitätsmessung
ein, enthält also die gleichen Gegenstände, wie später der Art. 4
seiner \textit{Allgemeinen Lehrsätze}. Anschliessend an diese Methoden müssen die
entsprechenden Sätze der Potentialtheorie schon frühe gefunden worden sein;
\clearpage\SourceLeaf{683}\ArithmeticHead{48}{L}{H. GEPPERT, ÜBER GAUSS’ ARBEITEN ZUR MECHANIK UND POTENTIALTHEORIE.}
\noindent dies gilt insbesondere a) von der \textsc{Poisson}schen Differentialgleichung, b) dem
ersten und zweiten Mittelwertsatz, c) der \textsc{Green}schen Formel (17).

1816 veröffentlicht \textsc{Gauss} die \textit{Demonstratio tertia} des Fundamentalsatzes
der Algebra, bei deren Analyse \textsc{Schlesinger}\footnote{\textsc{Schlesinger}, \textit{Über Gauss’ Arbeiten zur Funktionentheorie}, Werke X, 2, S. 162—165.} nachgewiesen hat, dass \textsc{Gauss}
die beiden Mittelwertsätze der Potentialtheorie in zwei Dimensionen verwandt
hat, die zur Umformung des angesetzten ebenen Doppelintegrals führen. In
der Tat beruht der zweite Mittelwertsatz ganz auf dem Satz von der scheinbaren
Grösse der \textit{Theoria attractionis}, dessen geometrische Einkleidung seine
Übertragung auf alle Dimensionen gestattet.

Eine weitere Bestätigung unserer Annahmen gibt das Werke III, S. 479—480
abgedruckte Nachlassstück, das auf etwa 1827 zu datieren ist und
schon den Zusammenhang zwischen Potential- und Funktionentheorie erkennen
lässt. Im ersten Stück [16] wird für zwei und drei Dimensionen die Konstanz
einer regulären Potentialfunktion aus der Konstanz auf dem Rande gefolgert,
womit ein erstes Ergebnis über die Wertverteilung gewonnen ist; als Hilfsmittel
dient eine (\textsc{Green}sche) Integralformel. Das zweite Stück [17] enthält
zunächst die Bemerkung, dass von einer Niveaufläche aus nach aussen bezw.
innen \(V\) monoton ab- bezw. zunimmt, den hieran naturgemäss anschliessenden
\textsc{Gauss}schen Satz (16) und schliesslich die Formel (17). Schon hier springt
der völlige Parallelismus zwischen der Theorie in zwei und drei Dimensionen
in die Augen\footnote{\textsc{Schlesinger} a.\,a.\,O. S. 186—187.}.

1829 schreibt \textsc{Gauss} die \textit{Theoria figurae fluidorum}\footnote{Werke V, S. 29—77, \textsc{Ostwalds} Klassiker, Band 135. Zur Kritik vergl. \textsc{Bolza}, a.\,a.\,O. S. 20—50.}, in der wieder Integralsätze
entwickelt werden, und zwar handelt es sich diesmal um die Umformung
eines ebenen Doppelintegrals in ein Linienintegral, also um den
\textsc{Green}schen Satz der Ebene in der von der heutigen etwas abweichenden
Form:
\[
\iint\limits_B\left\{\frac{\partial A}{\partial x}+\frac{\partial B}{\partial y}\right\}\,dx\,dy
=\int\limits_R(AY-BX)\,ds,\tag{19}
\]
\noindent worin \(R\) irgend eine geschlossene Raumkurve ist, deren Projektion die
Berandung von \(B\) bildet, und \(X,Y,Z\) die Richtungskosinus der positiven
\clearpage\SourceLeaf{684}\ArithmeticHead{49}{R}{IV. POTENTIALTHEORIE.}
\HistoricalFoot{X 2 Abh. 7.}{7}
\noindent Tangente an \(R\) sind. Die Formel wird durch die übliche Streifenzerlegung gewonnen;
sie umfasst bekanntlich den auch von \textsc{Gauss} frühzeitig erkannten
\textsc{Cauchy}schen Integralsatz.

Das Interesse an solchen Integralsätzen war um jene Zeit sehr rege.
Hatte auch die \textsc{Green}sche Abhandlung von 1828 keine sofortige Wirkung,
so verwandten doch schon 1833, also vor Veröffentlichung der \textit{Allgemeinen Lehrsätze},
\textsc{Duhamel}\footnote{\textsc{J. M. C. Duhamel}, \textit{Mémoire sur la méthode générale} .\, .\, . Journal de l’école polyt. 14, cah. 22,
1833, S. 20—77, bes. 68—69.} und \textsc{Lamé}\footnote{\textsc{G. Lamé}, \textit{Mémoire sur la propagation de la chaleur} .\, .\, ., ibidem S. 194—251, bes. S. 204.} beim Problem der Wärmeleitung die \textsc{Green}sche
Formel für drei Dimensionen. Doch hat \textsc{Gauss} als erster die Integration
durch Prismen- und Kegelzerlegung systematisch benutzt. \textsc{Beltrami}\footnote{\textsc{E. Beltrami}, \textit{Intorno ad alcuni punti della teoria del potenziale.} Mem. dell’Accad. di Bologna
9 (3), 1878, S. 451—475.} hat
später gezeigt, wie in diesen beiden Methoden, deren erstere als Spezialfall
der zweiten aufgefasst werden kann, die gesamten Formeln der Potentialtheorie,
insbesondere die \textsc{Green}schen Formeln, enthalten sind.

2. Die auf die konforme Abbildung und das logarithmische Potential
bezüglichen Studien von \textsc{Gauss} knüpfen an die Potentialgleichung an; schon
1802 beschäftigt sich \textsc{Gauss} mit der Gleichung
\[
\Delta u=\frac{\partial^2u}{\partial x^2}+\frac{\partial^2u}{\partial y^2}=0,
\]
\noindent deren Integral er zu \(u=f(x+iy)+\varphi(x-iy)\) angibt\footnote{\textsc{Schlesinger}, a.\,a.\,O. S. 156.}; hieran fügen sich unmittelbar
die oben behandelten Nachlassstücke aus Werke III, S. 479—480.
1822 erkennt er den Zusammenhang mit der Theorie der konformen Abbildung\footnote{\textsc{Schlesinger}, a.\,a.\,O. S. 173.}.
Das Zwei- und Dreidimensionale läuft in den \textsc{Gauss}schen Gedankengängen
parallel. Dies zeigt die Werke X, 1, S. 311 ff. abgedruckte,
etwa auf 1834 zu datierende Notiz, in der er sich mit der dem Abstände
umgekehrt proportionalen Anziehung in der Ebene beschäftigt, und die die
vollständige Beherrschung der Theorie des logarithmischen Potentials erkennen
lässt\footnote{\textsc{Laplace} hatte in seiner \textit{Mécanique céleste}, 2, cap. 13 zuerst bei der Untersuchung des Potentials eines unbegrenzten Zylinders das logarithmische Potential eingeführt.}. Dass beim Übergang von drei zu zwei Dimensionen im Ausdruck
des Potentials \(\frac1r\) durch \(\log\frac1r\) zu ersetzen sei, hatte \textsc{Gauss} schon früher
\clearpage\SourceLeaf{685}\ArithmeticHead{50}{L}{H. GEPPERT, ÜBER GAUSS’ ARBEITEN ZUR MECHANIK UND POTENTIALTHEORIE.}
\noindent bemerkt; in der Tat verwendet ja sein erstes algebraisches Theorem\footnote{Werke III, S. 112. \textsc{Schlesinger}, a.\,a.\,O. S. 189—191.} wesentlich
diese Tatsache, denn er formuliert diesen Satz so: Legt man in die
Wurzeln von \(f(x)=0\) gleiche Massen, so findet in den Wurzeln von \(f'(x)=0\)
Gleichgewicht statt. Da dieser Satz auf 1833—1836 zu datieren ist, hat sich
\textsc{Gauss} um jene Zeit offenbar sehr intensiv mit der Potentialtheorie beschäftigt,
was nicht zuletzt seinen Grund in den eifrig betriebenen erdmagnetischen
Untersuchungen findet.

Der Gipfelpunkt seiner Untersuchungen über das ebene Potential ist die
konforme Abbildung der Ellipse auf den Einheitskreis, die 1834—1839 durchgeführt
wurde und Werke X, 1, S. 311—320\footnote{\textsc{Schlesinger}, a.\,a.\,O. S. 192—194.} abgedruckt ist. Hier ist zunächst
bemerkenswert, dass \textsc{Gauss} die Existenz der Lösung des Randwertproblems,
die er in den Art. 29—37 seiner \textit{Allgemeinen Lehrsätze} entwickelt,
kennt, und gleichzeitig ist dies ein wichtiges Beispiel, das in einem speziellen
Falle die Methode der Lösung zeigt. Seine Methode ist geradezu die der \textsc{Green}schen
Funktion: Es wird die Masse \(2\pi A\) in den Mittelpunkt der Ellipse gelegt,
sie erzeuge das Potential \(V_1\), sodann wird diese Gesamtmasse \(2\pi A\) so
auf die Peripherie der Ellipse verteilt, dass ihr Potential \(V_2\) auf dieser Peripherie
mit \(V_1\) übereinstimmt. \(V_2-V_1\) ist dann die \textsc{Green}sche Funktion; ergänzt
man sie zu einer analytischen Funktion von \(x+iy\), so vermittelt sie bekanntlich,
als Exponent von \(e\) genommen, die konforme Abbildung der Ellipse auf den
Einheitskreis. Dies zeigt, dass \textsc{Gauss} erstens den Zusammenhang zwischen
dem Randwertproblem in der Ebene und der konformen Abbildung erkannt
hatte, sodann aber, dass er über das hinaus, was er in seiner Abhandlung niederlegte,
nicht nur den Existenz- und Eindeutigkeitssatz des ersten Randwertproblems,
sondern auch die Methode der Lösung mittels der \textsc{Green}schen
Funktion besass. Gleichzeitig erhellt, dass sich \textsc{Gauss} des völligen Parallelismus
der Theorie in zwei und drei Dimensionen bewusst war.

3. Über die Verbindung von Potentialtheorie und Erdmagnetismus bei
\textsc{Gauss} berichtet eingehend der Essay von \textsc{Schaefer}, Werke XI, 2, aus dem
wir nur die wichtigsten Datierungen zusammenstellen wollen. Die Wertverteilungssätze
der Art. 25—26 stammen spätestens aus dem Jahre 1831,
Den Fundamentalsatz des Randwertproblems, insbesondere in der von ihm
\clearpage\SourceLeaf{686}\ArithmeticHead{51}{R}{IV. POTENTIALTHEORIE.}
\HistoricalFoot{}{7\textsuperscript{*}}
\noindent verwandten Fassung, hat \textsc{Gauss} schon 1831 besessen, und offenbar auch den
zugehörigen Beweisgang, über dessen Schwierigkeit er sich im Briefe an \textsc{Encke}
vom 18. Aug. 1832\footnote{\textsc{Schaefer}, a.\,a.\,O. S. 101.} ausspricht; vorher wird er diesen Satz nicht besessen
haben, denn er enthält die Auflösung der begrifflichen Schwierigkeit, die in
den magnetischen Polen eines Körpers steckt. Mehrere bei \textsc{Schaefer}, S. 93
angeführte Briefstellen bestärken unsere Vermutung, dass \textsc{Gauss} den wesentlichen
Teil der Potentialtheorie schon in den ersten beiden Jahrzehnten des
neunzehnten Jahrhunderts besessen habe.

Schliesslich sei noch einiges über die bei \textsc{Gauss} auftretende Entwicklung
des Potentials nach Kugelfunktionen gesagt. Die Randwertaufgabe für den
Kreis kann nach der Manier von \textsc{Poisson}\footnote{\textsc{Poisson}, \textit{Mémoire sur la distribution de la chaleur dans les corps solides.} Journal de l’école
polyt. 12, cah. 19, 1823.} dadurch gelöst werden, dass die
Randbelegungsfunktion in eine trigonometrische Reihe entwickelt wird. Das
dreidimensionale Analogon hiervon führt \textsc{Gauss} durch, indem er die Belegung
einer Kugel nach Kugelfunktionen entwickelt, und damit die Entwicklung des
Kugelpotentials nach steigenden oder fallenden Potenzen des Abstandes vom
Kugelmittelpunkt gewinnt. Diese Methode hat \textsc{Gauss} in der \textit{Allgemeinen Theorie
des Erdmagnetismus} 1838\footnote{Werke V, S. 119—193.} durchgeführt; in den \textit{Allgemeinen Lehrsätzen} geht
er etwas weiter und löst im Art. 35 nach der gleichen Methode der Entwicklung
nach Kugelfunktionen auch das Randwertproblem für einen wenig von
der Kugel abweichenden Körper. Physikalisch können die einzelnen Terme
einer solchen Entwicklung als Potentiale von Dipolen, Tetrapolen usw. gedeutet
werden; auf diese Deutung bezieht sich ein Werke V, S. 631 abgedrucktes
Nachlassstück, dessen allgemeine Idee, die Kugelfunktionen durch
ihre Pole im Sinne der Potentialtheorie zu charakterisieren, erst später von
\textsc{Maxwell} vollständig durchgeführt worden ist\footnote{\textsc{E. W. Hobson}, \textit{The theory of spherical and ellipsoidal harmonics.} Cambridge 1931, S. 129—132.}.

Die Theorie der Kugelfunktionen gründet sich auf die vorausgegangenen
Arbeiten von \textsc{Laplace}\footnote{\textsc{Laplace}, \textit{Mécanique céleste}, 2, livre 3; 5, livre 11.} und \textsc{Legendre}\footnote{\textsc{Legendre}, \textit{Exercices de calcul intégral}, 1817, S. 247—273.}, die diese Funktionen zur Erledigung
des Attraktionsproblems der Ellipsoide einführten. \textsc{Gauss} gab den genannten
\clearpage\SourceLeaf{687}\ArithmeticHead{52}{L}{H. GEPPERT, ÜBER GAUSS’ ARBEITEN ZUR MECHANIK UND POTENTIALTHEORIE.}
\noindent Funktionen den Namen; er benutzte sie schon 1814 zur genäherten Quadratur.
Auf die Entwicklung eines homogenen Polynoms \(n\)-ten Grades \(f(x,y,z)\) auf
der Einheitskugel nach Kugelfunktionen bezieht sich ein weiteres \textsc{Gauss}sches
Nachlassstück in Werke V, S. 630—631; die Koeffizienten errechnen sich
leicht, wenn man von \(f\) sukzessive die \textsc{Laplace}sche \(\Delta\)-Operation bildet, die
nach \(\left[\frac n2\right]\) Schritten abbricht.
\clearpage\endgroup
